\section{Introduction}
\label{sec:introduction}

Environmental forecasting supports decisions that depend on both expected error and the timing or amplitude of changing conditions. A candidate forecast is therefore commonly compared with a meaningful reference forecast, rather than judged only by an isolated error value. Persistence-relative skill provides such a comparison for daily time series, while formal predictive-accuracy tests can assess whether an observed loss difference is statistically distinguishable from the reference \cite{hyndman2006,diebold1995}.

These error-based summaries are necessary but not exhaustive. A point forecast may reduce average squared error by attenuating fluctuations, particularly when the future is uncertain. Lower error can therefore coexist with a forecast trajectory whose variance or episode amplitude is smaller than that of the observations. Forecast verification distinguishes accuracy, association, calibration, and related properties because no single summary answers all of these questions \cite{gneiting2011,jolliffe2012}.

The distinction matters when evaluation results are converted into eligibility decisions. A cell that improves on persistence under RMSE may remain unsuitable for a decision that requires sensitivity to elevated pollution or preservation of dynamic spread. Variance and event diagnostics are established complements to error measures; they are not introduced here as new metrics. The unresolved practical issue is whether adding explicit fidelity requirements to an error-based eligibility rule materially changes the set of cells that pass.

The EMS audit motivating this study supports a deliberately bounded gap. Existing environmental forecasting studies commonly evaluate predictive error and, less consistently, dynamic-fidelity dimensions, while evidence remains limited on whether explicitly combining these dimensions in an eligibility rule materially changes decisions across sites, model families, and forecast horizons. This statement describes the audited corpus; it does not claim that no such study exists elsewhere.

We ask: \emph{Does incorporating dynamic-fidelity criteria alongside conventional error-based statistical evaluation materially alter model eligibility in multi-step daily PM$_{10}$ forecasting?} We answer with a post-evaluation audit of a fixed multi-station benchmark. The contributions are:
\begin{enumerate}
    \item an empirical audit of error--fidelity discordance;
    \item a characterisation of its variation across model families and horizons; and
    \item a quantification of how explicit fidelity requirements alter eligibility.
\end{enumerate}
